\section{Discussion}
\label{sec:discussion}

\subsection{The protocol changes the question, not only the scale}

The clearest methodological result is the 56-cell interpretation change. A raw
corrupted FP8--INT8 gap favors FP8 in 134 of 144 cells, but the clean treatment
already favors FP8 in every cell. After that baseline discrepancy is removed,
66 interactions are negative and the balanced mean is near zero. The adjustment
therefore prevents a clean-fidelity result from being relabelled as a
corruption-robustness result. This distinction is consequential even when the
same engine, images, and AP evaluator are used: the choice of contrast changes
what can legitimately be concluded.

The direct formulation also improves statistical efficiency and clarity. The
four treatment arms remain inside each image draw, retaining the covariance
induced by common images, bytes, preprocessing, and decoder. Subtracting two
independently summarized format effects would discard that pairing and obscure
which uncertainty sources are shared. The algebraic cancellation of the
FP32-typed reference further isolates the treatment comparison from the TF32
status of that historical reference, while the format-specific $Q$ and $E$
decompositions remain available for diagnosis.

\subsection{Clean fidelity, absolute accuracy, and interaction are distinct}

FP8's 1.40-AP matched-clean advantage is consistent across all 12 YOLO blocks.
That result is stronger and more uniform than the corruption interaction. The
exploratory grid mean is -0.02 AP with an interval crossing zero, whereas the
untouched VOC/KITTI estimate is -0.55 with an interval below zero. A negative
$\Delta E$ means that the FP8--INT8 gap contracts relative to matched clean; it
does not mean that INT8 retains more useful corrupted accuracy or that FP8 is
fragile in an absolute sense.

The absolute guardrail is therefore not optional. AP falls in 282 of 288
corrupted YOLO arms, and 35 arms are below 10 AP. The architecture stress cases
make the same point more sharply: the transferred RT-DETR INT8 recipe begins
with a 38.66-AP mean clean gap, then shows a negative interaction while
retaining only 1.97 mean corrupted AP. The interaction sign is arithmetically
correct but operationally unhelpful without the preceding fidelity and floor
checks. For deployment, the evidence should be read in order: matched-clean
accuracy, absolute corrupted accuracy, interaction, and only then latency or
engine size.

The term \emph{floor-guarded} is intentionally modest. We do not propose a
formal correction for bounded AP or a thresholded primary estimand. Instead, we
make failure floors visible and prevent the signed contrast from carrying a
claim it cannot support. Applications can add predeclared minimum-AP or
maximum-loss gates appropriate to their risk tolerance; such gates should be
specified before the interaction is inspected.

\subsection{Why the exploratory and holdout estimates differ}

The four-dataset landscape and the untouched VOC/KITTI rerun have different
conditional scopes. The broader grid includes COCO and TT100K, three capacities,
and evaluation partitions that were not uniformly independent of checkpoint
selection. It maps heterogeneity across a finite factorial design. The holdout
rerun removes evaluation-after-selection reuse for VOC and KITTI and retrains
six checkpoints under frozen split registries. Its -0.55-AP estimate is the
strongest split-independent evidence available here, but it does not cover COCO,
TT100K, additional training seeds, or alternative detector families.

The difference between -0.02 and -0.55 should therefore not be treated as a
failed replication. It shows that the interaction depends on the target regime.
This interpretation is consistent with the exploratory dataset margins,
leave-one-factor summaries, and 1.20-AP cell dispersion. Reporting one pooled
number without its finite-grid definition would hide that dependence. A future
confirmatory study should predeclare the deployment-weighted target, construct
untouched partitions for every domain, and propagate training, calibration,
engine-build, and acquisition uncertainty in one hierarchical design.

\subsection{Heterogeneity and object size}

The near-zero grid mean is produced by opposing condition effects rather than a
uniform absence of interaction. Dataset margins range from -0.46 on KITTI to
+0.61 on TT100K, and cell estimates span more than seven AP points. Alternative
weights and leave-one-factor summaries remain small but cross zero. This pattern
supports treatment- and regime-specific reporting rather than a single format
robustness rank.

The size endpoints also resist a simple narrative. The area-based
$\Delta\Psi$ interval reaches zero, whereas the TT100K height endpoint is
negative. These quantities compare the format interaction between small-like
and large-like objects; they do not measure absolute small-object resilience.
Size-specific AP is often low, particularly for TT100K XS instances, so floor
contraction can be stronger for the small endpoint. In addition, area thresholds
and original-box-height thresholds define different scientific objects and must
not be pooled. The results therefore do not support a universal claim that
quantization-corruption interaction is amplified on small objects.

\subsection{Treatment-level interpretation of INT8 and FP8}

The graph audit shows that both nominal formats are mixed-precision treatments
with architecture-dependent Q/DQ coverage. Calibration, scale granularity,
operator exclusions, TensorRT tactic selection, output parameterization, and
decoder behavior all contribute to the executable system. This treatment-level
view explains why the same nominal INT8 recipe can preserve YOLO accuracy yet
collapse on RT-DETR or RetinaNet. The portability stress test is valuable
precisely because it exposes failure rather than because it supplies a fair
cross-family leaderboard.

The evidence does not localize the cross-family INT8 collapse to a specific
layer, quantizer, runtime stage, or decoder. Repair would require output-level
parity checks, layer-wise sensitivity, alternative calibration, mixed-precision
rescue, quantizer-search methods, or quantization-aware training. Conversely,
FP8's close clean tracking across the six stress blocks is encouraging but
remains conditional on one recipe, one hardware/software stack, one training
seed, and one calibration list per dataset. Neither result establishes an
intrinsic property of the number format.

\subsection{Practical reporting recommendations}

For an executable detector comparison under corruption, we recommend the
following minimum report. First, define a matched-clean materialization and
show the clean format gap. Second, materialize each corruption independently of
precision and share the exact bytes across treatments. Third, estimate the
clean-adjusted format interaction inside common image resamples. Fourth, report
absolute corrupted AP or clean-to-corrupted loss, with any application-specific
accuracy gate declared in advance. Fifth, retain condition-level or stratified
results so that a balanced mean cannot conceal reversals. Finally, state the
full executable treatment and separate engine-only timing from end-to-end
latency.

These recommendations extend beyond INT8 and FP8. The same four-cell structure
can compare quantizers, mixed-precision policies, compilers, calibration
strategies, or hardware runtimes whenever a pre-existing clean discrepancy may
be confused with a distribution-shift interaction.

\subsection{Limitations}

Several limitations bound the conclusions. The corruption suite contains four
selected synthetic transformations and does not represent the distribution of
natural weather, sensor, compression, or domain shifts. Severity is ordinal and,
in the exploratory grid, stochastic realizations are not nested across levels.
The three-realization analysis is too small to estimate a population of possible
fields or angles.

The exploratory macro is a finite-grid arithmetic mean, not a deployment-
weighted population effect. The broad grid remains conditional on one primary
training seed and recorded checkpoint-selection workflows; the untouched rerun
covers only VOC and KITTI. The targeted seed crossing fixes YOLO11m, three
transfer datasets, and severity 5, and does not estimate full training or
calibration variance. Engine builds were not repeated, and the bootstrap
intervals condition on the frozen executable artifacts.

VOC and KITTI are evaluated through converted COCO-style annotations rather
than official challenge protocols. TT100K's sparse class universe can vary
under ordinary image resampling; the fixed-universe check covers only one stress
cell. The area and height endpoint families are not interchangeable. The
FP32-typed TensorRT reference may use TF32, although this does not affect the
direct $\Delta E$ algebra.

Runtime is measured on one RTX 5090/TensorRT stack and excludes preprocessing,
transfer, decoding, NMS, memory, power, and energy. Graph-level Q/DQ counts do
not reveal TensorRT kernel precision or fallback execution. The cross-family
extension uses point estimates, one training seed, one calibration identity,
and fixed recipes that were not optimized per architecture. Finally, the
compact submission package reproduces manuscript-level summaries from frozen
ledgers but does not redistribute datasets, checkpoints, engines, or raw
prediction payloads.
